% \section{Limitations}

% \paragraph{Static evaluation framework.}
% Our evaluation framework employs a static snapshot approach, assessing hallucinations at isolated decision points. While this method enhances reproducibility and attribution of results, it constrains our ability to determine whether incomplete or ambiguous actions will result in downstream errors. Consequently, certain borderline cases cannot be definitively categorized.

% \paragraph{Limitations of the LLM judge.}
% While our LLM-as-a-Judge framework achieves strong accuracy and consistency, it remains imperfect, with overall agreement below 90\%. Manual inspection reveals key limitations: (1) prompt design remains subjective and lacks formal validation, making it prone to bias; (2) the tradeoff between generalizable prompts and fine-grained evaluation limits precision, as case-specific rubrics require substantial manual effort; and (3) in rare cases, correct actions may be paired with unfaithful reasoning, exposing structural assumptions in our framework. Addressing these challenges will require more rigorous verification protocols and theoretical advances.

% \paragraph{Scope limited to digital LLM agents with textual observations.}
% Our benchmark exclusively targets digital LLM agents within generalist domains, employing purely textual observations. This narrow focus may limit the applicability of our findings to embodied agents, multi-modal environments, or highly specialized tasks, as the operational contexts for these settings differ substantially.

% \paragraph{Baseline agent constraints.}
% Our evaluation relies solely on baseline agents provided within each specific environment, utilizing a ReAct-style generation pnrocess for decision-making and action-taking. The scope of hallucination examined in our study is confined to the working memory available to baseline agents, excluding any retrieved information from long-term memory stores such as semantic, episodic, or procedural memory \cite{sumers2023cognitive}. Additionally, our framework does not address the learning behaviors of language agents. As future research develops more sophisticated LLM agents featuring these advanced capabilities, the identified risk settings and taxonomy of hallucinations might correspondingly evolve.